\subsection*{Reading beyond the last two glyphs}
Can a pictured book prefix continue to influence a writer after its first error? We reserved two entire classical Chinese works, the \emph{Analects} and \emph{Chu Ci}, before fitting a causal image reader from scratch. Given 32 printed glyphs, the reader could predict a training bitmap, copy a visible cell or use counts of short preceding sequences. With the original history supplied, this mixture recovered 2,744 of 28,880 next glyphs. The matched image reader without counts recovered 2,613 and the count predictor recovered 2,354. Paired differences were 0.45 percentage points (95\% volume-bootstrap interval 0.12 to 0.78) and 1.35 points (0.80 to 2.03). A byte-input mixture recovered 2,809. Its difference from the visual mixture was unresolved across the 29 held volumes (Fig.~\ref{fig:book-prefix}).

With self-feedback, the visual mixture matched 839 positions, falling to 426 after exchanging the first 30 pictured glyphs with a prefix from the other work. The original prefix helped in both books. An exchange can change subject, diction and verse rhythm together. The unassisted neural reader recovered 969 positions during free writing, reversing the ordering observed with correct history. All 1,805 reference continuations were absent as complete sixteen-cell sequences from training. No writer recovered even the first eight glyphs of a complete reference. The two initial examples show why a legible output can still be a poor continuation. Both cycle through a few glyphs from the prompt. A single historical continuation provides a literal recovery test; alternate readings require separate assessment.
